\section{Programmazione dinamica}
\subsection{Memoizzazione}
\subsubsection*{Concetto di memoizzazione per istanze ripetute}
Il paradigma della programmazione dinamica è un caso particolare di divide and conquer. Viene utilizzato nei problemi in cui
alcune sottoistanze vengono risolte più volte in punti diversi dell'albero della ricorsione. Per evitare il lavoro ridondante
e migliorare l'efficienza dell'algoritmo, la programmazione dinamica sfrutta la memoizzazione che consiste nel tenere traccia
delle varie sottoistanze risolte con le rispettive soluzioni, in modo da non doverle ricalcolare quando si presentano nuovamente.

La presenza di sottoistanze uguali è abbastanza rara, per cui il paradigma di programmazione dinamica viene impiegato
raramente e solo in determinati problemi con strutture particolari.

\subsubsection*{Approccio bottom-up iterativo e ordine di visita delle sottoistanze}
In alternativa all'approccio top-down ricorsivo tipicamente usato nel D\&C, è possibile risolvere i problemi di programmazione
dinamica con un approccio bottom-up iterativo. Durante questo processo, si esplorano tutte le sottoistanze dal caso base fino
a raggiungere l'istanza principale richiesta dall'utente e si memorizzano le soluzioni temporanee in una struttura dati, come
ad esempio una tabella.

Per fare ciò, in fase di progettazione, è necessario identificare un ordine di visita delle sottoistanze in modo tale da
garantire che una sottoistanza venga risolta solo dopo aver risolto tutte le sottoistanze da cui essa dipende. In altre
parole è necessario individuare un ordinamento topologico delle sottoistanze.

\subsection{Calcolo dell'\texorpdfstring{\(n\)}{n}-esimo numero di Fibonacci}
\subsubsection*{Formalizzazione del problema}
L'\(n\)-esimo numero di Fibonacci è definito ricorsivamente come segue:
\[F(n) = \begin{cases}
	n & \text{se } n = 0,1 \\
	F(n-1) + F(n-2) & \text{se } n > 1
\end{cases} \qquad \text{con} \qquad \mathcal{I} = \mathbb{Z}_{+,0} \qquad \mathcal{S} = \mathbb{Z}_{+,0}\]
L'istanza del problema è data da un numero \(n\), per cui la taglia risulta essere \(\log_2(n)\).

\subsubsection*{Soluzione ricorsiva non efficiente}
Applicando puramente la definizione ricorsiva illustrata sopra si ottiene il seguente algoritmo.

\begin{algo}{REC\_FIB}{n}{intero positivo $n \in \mathbb{Z}_{+,0}$}{$n$-esimo numero di Fibonacci $F(n) \in \mathbb{Z}_{+}$}
	\If{\(n = 0\) \textbf{or} \(n = 1\)}
		\State \Return \(n\)
	\Else
		\State \Return \(\text{REC\_FIB}(n-1) + \text{REC\_FIB}(n-2)\)
	\EndIf
\end{algo}

\vspace{5pt}

\noindent
L'equazione alle ricorrenze dell'algoritmo è la seguente:
\[T(n) = \begin{cases}
	0 & \text{se } n = 0,1 \\
	T(n-1) + T(n-2) + 1 & \text{se } n > 1
\end{cases}\]
Osservando che \(T(n-1) \geq T(n-2)\), si calcola un lower bound della complessità temporale che risulta essere doppiamente
esponenziale nella taglia dell'istanza, per cui fortemente inefficiente.
\begin{align*}
	T(n) \;\; &\geq \;\; 2T(n-2) + 1 \\[3pt]
	&= \;\; \Omega \! \left( 2^{n/2} \right) \; = \; \Omega \! \left( \sqrt{2}^n \right) \; = \; \Omega \! \left( 2^{2^{(\log_2 n) - 1}} \right) \; = \; \Omega \! \left( 2^{2^{\text{taglia(n)} - 1}} \right)
\end{align*}

\newpage

\subsubsection*{Soluzione iterativa con memoizzazione}
Siccome l'\(n\)-esimo numero di Fibonacci dipende solo dai due numeri precedenti, si osserva che è sufficiente tenere in
memoria i due numeri di Fibonacci più recenti appena calcolati per poter calcolare il successivo. Questo tipo di approccio
è anche detto bottom-up.

\begin{algo}{IT\_FIB}{n}{intero positivo $n \in \mathbb{Z}_{+,0}$}{$n$-esimo numero di Fibonacci $F(n) \in \mathbb{Z}_{+}$}
	\If{\(n = 0\) \textbf{or} \(n = 1\)}
		\State \Return \(n\)
	\EndIf

	\State
	\State \(F[0] \gets 0\)
	\State \(F[1] \gets 1\)
	\For{\(i \gets 2\) \textbf{to} \(n\)}
		\State \(F[i] \gets F[i-1] + F[i-2]\)
	\EndFor
	\State \Return \(F[n]\)
\end{algo}

~

\noindent
Siccome l'algoritmo non è ricorsivo, la complessità risulta più facile da calcolare e asintoticamente è esponenziale nella
taglia dell'istanza, per cui comunque poco efficiente.
\[T(n) = n + 1 = 2^{\log_2 n} - 1 = 2^\text{taglia(n)} - 1\]


\subsubsection*{Soluzione analitica con formula di Binet ed esponenziazione di un numero}
Riscrivendo la definizione ricorsiva di Fibonacci in forma matriciale e calcolando successivamente gli autovalori e autovettori
della matrice di trasferimento è possible convertire la definizione ricorsiva in una formula chiusa nota come formula di Binet.
\[\begin{pmatrix} F_n \\ F_{n-1} \end{pmatrix} = \begin{pmatrix} 1 & 1 \\ 1 & 0 \end{pmatrix} \begin{pmatrix} F_{n-1} \\ F_{n-2} \end{pmatrix} \qquad \rightarrow \qquad \begin{pmatrix} F_n \\ F_{n-1} \end{pmatrix} = \begin{pmatrix} 1 & 1 \\ 1 & 0 \end{pmatrix}^{n-1} \begin{pmatrix} F_{1} \\ F_{0} \end{pmatrix} = \begin{pmatrix} 1 & 1 \\ 1 & 0 \end{pmatrix}^{n-1} \begin{pmatrix} 1 \\ 0 \end{pmatrix}\]
\vspace{-7pt}
\[\text{autovalori ed autovettori} : \quad \lambda_1 = \frac{1 + \sqrt{5}}{2}, \quad \lambda_2 = \frac{1 - \sqrt{5}}{2}, \quad v_1 = \left(\frac{1 + \sqrt{5}}{2}, 1\right), \quad v_2 = \left(\frac{1 - \sqrt{5}}{2}, 1\right)\]
\vspace{-14pt}
\begin{align*}
	\begin{pmatrix} 1 & 1 \\ 1 & 0 \end{pmatrix}^{n-1} &= \begin{pmatrix} {v_1}^T \\ {v_2}^T \end{pmatrix} \cdot \begin{pmatrix} \lambda_1 & 0 \\ 0 & \lambda_2 \end{pmatrix}^{n-1} \cdot \begin{pmatrix} {v_1}^T & {v_2}^T \end{pmatrix}^{-1} = \\
	&= \begin{pmatrix} \frac{1 + \sqrt{5}}{2} & \frac{1 - \sqrt{5}}{2} \\[8pt] 1 & 1 \end{pmatrix} \cdot \begin{pmatrix} \left(\frac{1 + \sqrt{5}}{2}\right)^{n-1} & 0 \\[-2pt] 0 & \left(\frac{1 - \sqrt{5}}{2}\right)^{n-1} \end{pmatrix} \cdot \begin{pmatrix} \frac{1}{\sqrt{5}} & -\frac{1 - \sqrt{5}}{2\sqrt{5}} \\[8pt] -\frac{1}{\sqrt{5}} & \frac{1 + \sqrt{5}}{2\sqrt{5}} \end{pmatrix} = \\
	&= \begin{pmatrix} \frac{1}{\sqrt{5}} \cdot \left[ \left(\frac{1 + \sqrt{5}}{2}\right)^n - \left(\frac{1 - \sqrt{5}}{2}\right)^n \right] & \dots \\[8pt] \frac{1}{\sqrt{5}} \cdot \left[ \left(\frac{1 + \sqrt{5}}{2}\right)^{n-1} - \left(\frac{1 - \sqrt{5}}{2}\right)^{n-1} \right] & \dots \end{pmatrix} = \begin{pmatrix} \frac{1}{\sqrt{5}} \cdot \left( {\lambda_1}^n - {\lambda_2}^n \right) & \dots \\[8pt] \frac{1}{\sqrt{5}} \cdot \left( {\lambda_1}^{n-1} - {\lambda_2}^{n-1} \right) & \dots \end{pmatrix}
\end{align*}
\[\begin{pmatrix} F_n \\ F_{n-1} \end{pmatrix} = \begin{pmatrix} 1 & 1 \\ 1 & 0 \end{pmatrix}^{n-1} \begin{pmatrix} 1 \\ 0 \end{pmatrix} = \begin{pmatrix} \frac{1}{\sqrt{5}} \cdot \left( {\lambda_1}^n - {\lambda_2}^n \right) & \dots \\[8pt] \frac{1}{\sqrt{5}} \cdot \left( {\lambda_1}^{n-1} - {\lambda_2}^{n-1} \right) & \dots \end{pmatrix} \cdot \begin{pmatrix} 1 \\ 0 \end{pmatrix} \quad \rightarrow \quad F(n) = \frac{1}{\sqrt{5}} \cdot \left( {\lambda_1}^n - {\lambda_2}^n \right)\]
~

\noindent
La complessità della formula analitica di Binet è dominata dal calcolo delle potenze di numeri reali, che risultano
essere lineari nella taglia dell'istanza dell'esponente per come visto nell'algoritmo efficiente di esponenziazione,
per cui \(T(n) = \Theta(\log_2 n)\).

Siccome si lavora con numeri reali, però, servono analisi più approfondite per valutare sia la complessità, siccome
la taglia di un numero reale dipende dalla rappresentazione in virgola mobile utilizzata, sia l'accuratezza dei risultati,
siccome le rappresentazioni in virgola mobile introducono errori di approssimazione che non possono essere trascurati.
Per questo motivo l'algoritmo basato sulla formula di Binet si accenna solo dal punto di vista teorico.

\subsubsection*{Soluzione ricorsiva con memoizzazione}
È possibile sviluppare un approccio ricorsivo con memoizzazione, in cui la funzione ricorsiva costruisce e memorizza le
soluzioni intermedie in un vettore come nell'algoritmo iterativo, piuttosto che costruirle da zero ogni volta e dimenticarle.
La complessità di questo approccio è uguale a quella dell'algoritmo iterativo, ovvero esponenziale nella taglia dell'istanza.

L'implementazione è composta da una funzione di inizializzazione INIT\_FIB(\(n\)) che costruisce e inizializza il vettore
dove salvare le soluzioni intermedie e da una funzione ricorsiva REC\_FIB\_MEM(\(n\)) che costruisce le soluzioni intermedie
in modo ricorsivo memorizzandole nel vettore.

\begin{algo}{INIT\_FIB}{n}{intero positivo $n \in \mathbb{Z}_{+,0}$}{$n$-esimo numero di Fibonacci $F(n) \in \mathbb{Z}_{+}$}
	\If{\(n = 0\) \textbf{or} \(n = 1\)}
		\State \Return \(n\)
	\EndIf

	\State
	\State \(F[0] \gets 0\)
	\State \(F[1] \gets 1\)
	\For{\(i \gets 2\) \textbf{to} \(n\)}
		\State \(F[i] \gets 0\)
	\EndFor
	\State \Return \text{REC\_FIB\_MEM}\(n\)
\end{algo}

\begin{algo}{REC\_FIB\_MEM}{n}{intero positivo $n \in \mathbb{Z}_{+,0}$}{$n$-esimo numero di Fibonacci $F(n) \in \mathbb{Z}_{+}$}
	\If{\(F[n] = 0\)}
		\State \(F[n] \gets \text{REC\_FIB\_MEM}(n-1) + \text{REC\_FIB\_MEM}(n-2)\)
	\EndIf
	\State \Return \(F[n]\)
\end{algo}

\newpage

\subsection{Longest Common Subsequence - LCS}
\subsubsection*{Definizione di sottosequenza e di sottostringa di una stringa}
La sottosequenza di una stringa \(X\) è una successione di simboli non necessariamente contigui che mantengono lo stesso
ordinamento con cui si trovano nella stringa \(X\). Una sottostringa è una sottosequenza particolare in cui è richiesto
che i simboli devono essere contigui.
\[\begin{array}{c c}
	\displaystyle \binom{n}{k} = \frac{n!}{k! \, (n-k)!} & \quad \text{\# di sottosequenze di lunghezza \(k\) di una stringa di lunghezza \(n\)} \\[13pt]
	\displaystyle \sum_{k=0}^{n} \binom{n}{k} = 2^n & \quad \text{\# di tutte le sottosequenze di una stringa di lunghezza \(n\)}
\end{array}\]

\noindent
Si nota che per ottenere una sottosequenza di lunghezza \(k\), si scelgono a caso \(k\) simboli dalla stringa \(X\) di lunghezza
\(n\) senza tenere conto dell'ordine. È equivalente a cercare un sottoinsieme di \(k\) elementi da un insieme di \(n\)
elementi, per questo motivo si utilizza il coefficiente binomiale. L'ordine dei simboli nella sottosequenza è imposto a
posteriori della scelta, in base all'ordine con cui si trovano nella stringa \(X\).

\subsubsection*{Formalizzazione del problema}
Date due stringhe \(X\) e \(Y\) rispettivamente di lunghezza \(m\) e \(n\), si vuole calcolare la sottosequenza più lunga
\(Z\) che sia comune ad entrambe le stringhe \(X\) e \(Y\). La taglia dell'istanza è data da \(m + n\).

\subsubsection*{Descrizione dell'algoritmo LCS(\(X\), \(Y\)) ricorsivo con memoizzazione}
\begin{itemize}
	\item \textbf{caso base}: \\
	se una delle due stringhe \(X\) o \(Y\) è vuota, la sottosequenza comune più lunga è anch'essa vuota
	\item \textbf{passo ricorsivo - caso 1}: \\
	se le due stringhe \(X = \; \langle X', a \rangle\) e \(Y = \; \langle Y', a \rangle\) terminano con lo stesso simbolo \(a\), allora la sottosequenza
	comune più lunga è data dalla sottosequenza comune più lunga tra \(X'\) e \(Y'\) (\(\text{LCS}(X', Y')\)) con l'aggiunta del
	simbolo \(a\) alla fine
	\item \textbf{passo ricorsivo - caso 2}: \\
	se le due stringhe \(X = \; \langle X', a \rangle\) e \(Y = \; \langle Y', b \rangle\) terminano con simboli diversi \(a\) e \(b\), allora la sottosequenza
	comune più lunga sarà data da quella di lunghezza massima tra la sottosequenza comune più lunga tra \(X'\) e \(Y\) (\(\text{LCS}(X', Y)\)) e
	quella tra \(X\) e \(Y'\) (\(\text{LCS}(X, Y')\))
\end{itemize}
\vspace{5pt}

\noindent
Per identificare i prefissi di \(X\) e \(Y\) da passare alle chiamate ricorsive, si utilizzano rispettivamente gli
indici \(i\) e \(j\) che puntano all'ultimo simbolo di \(X\) e \(Y\) da considerare. Di seguito sono riportate la proprietà
di sottostruttura e il calcolo della lunghezza \(l\) della LCS indicate come equazioni alle ricorrenze.
\[Z(X_{0:i}, Y_{0:j}) \; = \text{LCS}(i, j) \; = \; \begin{cases}
	\; \varepsilon & \text{caso base} - \text{per } i = 0 \vee j = 0 \\[3pt]
	\; \text{LCS}(i-1, j-1) \;\; \cup \; \langle X_i \rangle & \text{caso 1} - \text{se } X_i = Y_j \\[3pt]
	\; \max \{\; \text{LCS}(i-1, j), \; \text{LCS}(i, j-1)\; \} & \text{caso 2} - \text{se } X_i \neq Y_j
\end{cases}\]
\[\text{len}(Z(X_{0:i}, Y_{0:j})) \; = \; l(i, j) \; = \; \begin{cases}
	0 & \text{caso base} - \text{per } i = 0 \vee j = 0 \\[3pt]
	l(i-1, j-1) + 1 & \text{caso 1} - \text{se }X_i = Y_j \\[3pt]
	\max \{ l(i-1, j), l(i, j-1) \} & \text{caso 2} - \text{se } X_i \neq Y_j
\end{cases}\]

\vspace{5pt}

\noindent
L'algoritmo viene definito con approccio bottom-up, ovvero si parte da un caso base semplice \(i = j = 0\) e si
espandono le soluzioni intermedie facendo variare progressivamente un indice alla volta fino ad arrivare alla soluzione
finale \(i = m\) e \(j = n\). Ogni soluzione intermedia viene memorizzata in una tabella bidimensionale, come illustrato
successivamente.

\newpage

\subsubsection*{Tabella dei risultati (struttura dati di memoizzazione)}
Le soluzioni intermedie vengono memorizzate in una tabella bidimensionale di dimensione \((m+1) \times (n+1)\) con le
seguenti caratteristiche:
\begin{itemize}
	\item la cella \((i,j)\) fa riferimento alla chiamata \(LCS(i, j)\)
	\item il valore memorizzato nella cella \((i,j)\) è la lunghezza \(l(i,j)\) della \(LCS(i, j)\)
	\item la soluzione finale \(LCS(X, Y)\) è contenuta nella cella \((m,n)\) (in basso a destra) della tabella
	\item in ogni cella è riportata una freccia che indica la cella riferita alla chiamata ricorsiva effettuata per
	calcolare il valore della cella stessa, con le seguenti regole:
	\begin{itemize}[topsep=-2pt]
		\item se \(X_i = Y_j\) la freccia è diagonale e punta alla cella \((i\!-\!1, \; j\!-\!1)\)
		\item se \(X_i \neq Y_j\) la freccia è orizzontale o verticale e punta rispettivamente alla cella \((i\!-\!1, \; j)\)
		o alla cella \((i, \; j\!-\!1)\) a seconda di quale delle due chiamate ricorsive ha restituito il valore più grande
		(alcune volte è indipendente scegliere tra orizzontale o verticale se i valori sono uguali)
	\end{itemize}
	\item per l'approccio di costruzione bottom-up, il popolamento della tabella parte dal caso base \((0,0)\) (in alto
	a sinistra) e prosegue per righe (approccio row-major) oppure per colonne (approccio column-major), in base a che
	indice si decide di far variare per primo
\end{itemize}

\noindent
Di seguito ne è riportato un esempio per le stringhe \(X = \, \langle A, B, C, B, B, D \rangle\) e \(Y = \; \langle A, D, C, C, B, D \rangle\).
\begin{center}
	\renewcommand{\arraystretch}{1.3}
	\setlength{\tabcolsep}{12pt}
	\begin{tabular}{c|c|c|c|c|c|c|c|}
	\diagbox{x}{y} & \(\begin{array}{c} \!\!\!\!Y_0\!\!\!\! \\ \varepsilon \end{array}\) & \(\begin{array}{c} \!\!\!\!Y_1\!\!\!\! \\ A \end{array}\) & \(\begin{array}{c} \!\!\!\!Y_2\!\!\!\! \\ D \end{array}\) & \(\begin{array}{c} \!\!\!\!Y_3\!\!\!\! \\ C \end{array}\) & \(\begin{array}{c} \!\!\!\!Y_4\!\!\!\! \\ C \end{array}\) & \(\begin{array}{c} \!\!\!\!Y_5\!\!\!\! \\ B \end{array}\) & \(\begin{array}{c} \!\!\!\!Y_6\!\!\!\! \\ D \end{array}\) \\
	\hline
	\(X_0: ~ \varepsilon\) & 0 & 0 & 0 & 0 & 0 & 0 & 0 \\
	\hline
	\(X_1: ~ A\) & 0 & $\nwarrow$ 1 & $\leftarrow$ 1 & $\leftarrow$ 1 & $\leftarrow$ 1 & 1 & 1 \\
	\hline
	\(X_2: ~ B\) & 0 & 1 & 1 & 1 & $\uparrow$ 1 & 2 & 2 \\
	\hline
	\(X_3: ~ C\) & 0 & 1 & 1 & 2 & $\nwarrow$ 2 & 2 & 2 \\
	\hline
	\(X_4: ~ B\) & 0 & 1 & 1 & 2 & $\uparrow$ 2 & 3 & 3 \\
	\hline
	\(X_5: ~ B\) & 0 & 1 & 1 & 2 & 2 & $\nwarrow$ 3 & 3 \\
	\hline
	\(X_6: ~ D\) & 0 & 1 & 2 & 2 & 2 & 3 & $\nwarrow$ 4 \\
	\hline
	\end{tabular}
\end{center}

\subsubsection*{Ricostruzione della LCS a partire dalle frecce della tabella}
La ricostruzione della sottosequenza comune più lunga \(Z\) equivale ad una visita in post-order dell'ordinamento topologico
delle frecce contenute nella tabella \(\underline{b}\). La prima chiamata deve essere effettuata con gli indici \(i = m\)
e \(j = n\) che puntano all'ultima cella della tabella. La visita termina quando si raggiunge una cella con indice \(i = 0\)
o \(j = 0\) (caso base).

\begin{algo}{PRINT\_LCS}{$\underline{X}$, $\underline{b}$, $i$, $j$}{stringa $\underline{X}$ per recupero elementi, tabella $\underline{b}$ con le frecce e indici $i$, $j$}{nessun return, stampa la LCS($X$, $Y$) a standard output}
	\If{\(m = 0\) \textbf{or} \(n = 0\)} \Comment{invocazione sul caso base \(\rightarrow\) nessuna azione necessaria}
		\State \Return \(-\)
	\EndIf
	\If{\(b[i, j] = \text{``}\nwarrow\text{''}\)} \Comment{caso 1 - cella diagonale}
		\State \text{PRINT\_LCS}(\(\underline{X}, \underline{b}, i-1, j-1\))
		\State print(\(X_i\))
	\ElsIf{\(b[i, j] = \text{``}\uparrow\text{''}\)} \Comment{caso 2 - cella verticale}
		\State \text{PRINT\_LCS}(\(\underline{X}, \underline{b}, i-1, j\))
	\Else \(\quad (\text{ovvero per }b[i, j] = \text{``}\leftarrow\text{''}\)) \Comment{caso 2 - cella orizzontale}
		\State \text{PRINT\_LCS}(\(\underline{X}, \underline{b}, i, j-1\))
	\EndIf
\end{algo}

\newpage

\subsubsection*{Implementazione iterativa bottom-up}
\vspace{-7pt}
\begin{algo}{IT\_LCS}{$\underline{X}$, $\underline{Y}$}{stringhe $\underline{X}$ e $\underline{Y}$ rispettivamente di lunghezza $m$ e $n$}{$(l(m, n), \underline{b})$ lunghezza e tabella con le frecce per ricostruire la LCS($X$, $Y$)}
	\State \(m \gets \text{len}(X)\)
	\State \(n \gets \text{len}(Y)\)
	\For {\(i \gets 0\)}{\(m\)} \(l[i, 0] \gets 0\) \Comment{inizializzazione caso base della prima riga} \EndFor
	\For {\(j \gets 0\)}{\(n\)} \(l[0, j] \gets 0\) \Comment{inizializzazione caso base della prima colonna} \EndFor

	\State
	\For{\(i \gets 1\)}{\(m\)} \Comment{popolamento della tabella row-major}
		\For{\(j \gets 1\)}{\(n\)}
			\If{\(X_i = Y_j\)} \Comment{passo ricorsivo - caso 1}
				\State \(l[i, j] \gets l[i-1, j-1] + 1 \quad-\quad b[i, j] \gets \text{``}\nwarrow\text{''}\)
			\ElsIf{\(l[i-1, j] \geq l[i, j-1]\)} \Comment{passo ricorsivo - caso 2 - cella verticale}
				\State \(l[i, j] \gets l[i-1, j] \quad-\quad b[i, j] \gets \text{``}\uparrow\text{''}\)
			\ElsIf{\(l[i-1, j] < l[i, j-1]\)} \Comment{passo ricorsivo - caso 2 - cella orizzontale}
				\State \(l[i, j] \gets l[i, j-1] \quad-\quad b[i, j] \gets \text{``}\leftarrow\text{''}\)
				\State \(b[i, j] \gets \text{``}\leftarrow\text{''}\)
			\EndIf
		\EndFor
	\EndFor
	\State \Return (\(l[m][n], \;\; \underline{b}\))
\end{algo}

\subsubsection*{Implementazione ricorsiva top-down con memoizzazione}
\vspace{-7pt}
\begin{algo}{INIT\_LCS}{$\underline{X}$, $\underline{Y}$}{stringhe $\underline{X}$ e $\underline{Y}$ rispettivamente di lunghezza $m$ e $n$}{$(l(m, n), \underline{b})$ lunghezza e tabella con le frecce per ricostruire la LCS($X$, $Y$)}
	\State \(m \gets \text{len}(X)\)
	\State \(n \gets \text{len}(Y)\)
	\If{\(m = 0\) \textbf{or} \(n = 0\)} \Comment{invocazione sul caso base \(\rightarrow\) nessuna azione necessaria}
		\State \Return \(0\)
	\EndIf
	\For {\(i \gets 0\)}{\(m\)} \(l[i, 0] \gets 0\) \Comment{inizializzazione caso base della prima riga} \EndFor
	\For {\(j \gets 0\)}{\(n\)} \(l[0, j] \gets 0\) \Comment{inizializzazione caso base della prima colonna} \EndFor
	\For {\(i \gets 1\)}{\(m\)} \Comment{inizializzazione con valore sentinella \(-1\)}
		\For {\(j \gets 1\)}{\(n\)} \Comment{per identificare i casi non ancora gestiti}
			\State \(l[i,j] \gets -1\)
		\EndFor
	\EndFor
	\State \Return \text{REC\_LCS}\((\underline{X}, \underline{Y},m, n)\)
\end{algo}

\begin{algo}{REC\_LCS}{$\underline{X}$, $\underline{Y}$, $m$, $n$}{stringhe $\underline{X}$ e $\underline{Y}$ e rispettive lunghezze $m$ e $n$}{$(l(m, n), \underline{b})$ lunghezza e tabella con le frecce per ricostruire la LCS($X$, $Y$)}
	\If{\(l[m, n] = -1\)} \Comment{caso non ancora gestito}
		\If{\(X_i = Y_j\)} \Comment{passo ricorsivo - caso 1}
			\State \(l[i, j] \gets 1 + \text{REC\_LCS}(\underline{X}, \underline{Y}, i-1, j-1) \quad-\quad b[i, j] \gets \text{``}\nwarrow\text{''}\)
		\ElsIf{\(\text{REC\_LCS}(\underline{X}, \underline{Y}, i\!-\!1, j) \geq \text{REC\_LCS}(\underline{X}, \underline{Y}, i, j\!-\!1)\)} \Comment{caso 2 - cella verticale}
			\State \(l[i, j] \gets l[i-1, j] \quad-\quad b[i, j] \gets \text{``}\uparrow\text{''}\)
		\Else \(\quad (\text{ovvero per }l[i-1, j] < l[i, j-1])\) \Comment{caso 2 - cella orizzontale}
			\State \(l[i, j] \gets l[i, j-1] \quad-\quad b[i, j] \gets \text{``}\leftarrow\text{''}\)
		\EndIf
	\EndIf
	\State \Return (\(l[m][n], \;\; \underline{b}\))
\end{algo}

\newpage

\subsubsection*{Complessità temporale}
La complessità temporale dell'algoritmo LCS è \(O(m \cdot n)\) sia per l'implementazione iterativa che per quella ricorsiva
con memoizzazione, per come di seguito spiegato:
\begin{itemize}
	\item nella versione iterativa, si visitano tutte le \(n \cdot m\) celle della tabella e per ogni cella viene fatto
	un numero costante di operazioni
	\item nella versione ricorsiva con memoizzazione, ogni cella della tabella deve essere inizializzata per cui il
	ragionamento è lo stesso della versione iterativa
\end{itemize}
Tale complessità, seppur quadratica per \(m = n\), è comunque molto più efficiente rispetto alla complessità esponenziale
dell'algoritmo banale brute-force che esplora tutte le possibili sottosequenze di \(X\) e \(Y\) per identificare quella
in comune più lunga.

Con alcuni accorgimenti è possibile ridurre la complessità temporale fino a \(O(m \cdot n / \log_2 n)\).

\subsubsection*{NO-SETH e Fine Grained Complexity}
Il problema LCS è definito come NO-SETH, ovvero è un algoritmo che se dovesse essere risolto in tempo polinomiale subquadratico
\(O(n^{2-\varepsilon})\), allora sarebbe possibile risolvere in tempo polinomiale molti altri problemi che attualmente hanno
solo soluzioni con complessità esponenziale. Questo studio è svolto nell'ambito della ``Fine Grained Complexity'', ovvero
una branca dello studio della complessità che si occupa di identificare i valori precisi delle complessità di algoritmi per
studiare come piccoli miglioramenti su alcuni algoritmi possano portare al passaggio da esponenziale a polinomiale per molti
altri problemi.

\subsubsection*{Complessità spaziale}
La complessità spaziale dell'algoritmo LCS è \(O(m \cdot n)\), ovvero è proporzionale al numero di celle della tabella in
cui vengono salvati i risultati. Corrisponde a tutte le possibili combinazioni di prefissi di \(X\) e \(Y\) che possono
essere passati alle chiamate ricorsive.

Se si è interessati a conoscere solo la lunghezza della LCS, è possibile ridurre la complessità spaziale a \(O(2 \cdot \min\{m, n\})\)
memorizzando solo le due righe (o colonne) più recenti della tabella (quella precedente e quella che si sta popolando),
in quando le equazioni alle ricorrenze dipendono solo dai valori delle celle adiacenti (diagonale, orizzontale e verticale).

\subsubsection*{Osservazioni sulla struttura dati per la memoizzazione}
Per gestire la memoizzazione è richiesta una struttura dati con le seguenti operazioni:
\begin{itemize}
	\item ricerca tramite la sottoistanza (coppia di indici (\(i, j\))) che restituisce la soluzione della sottoistanza cercata
	o \(-1\) se tale sottoistanza non è ancora stata gestita
	\item inserimento delle soluzioni per le sottoistanze risolte
\end{itemize}
La struttura più semplice che garantisce tali operazioni con complessità \(O(1)\) è una matrice bidimensionale, ma ha lo
svantaggio che lo spazio occupato è fisso di \(O(m \cdot n)\) indipendentemente dal numero di sottoistanze effettivamente
gestite.

Un'alternativa più efficiente a livello spaziale sono il dizionario o mappa (eventualmente semplificati in quanto non è
richiesta la cancellazione) in quanto occupano lo spazio necessario per memorizzare solo le sottoistanze effettivamente
gestite. Lo svantaggio è che hanno complessità di ricerca e inserimento di \(O(\log_2 (m \cdot n))\). Risultano, quindi,
più efficienti solo se il numero di sottoistanze effettivamente gestite è minore di \((m \cdot n)/\log_2 (m \cdot n)\).

\newpage

\subsection{Matrix Chain Multiplication - MCM}
\subsubsection*{Introduzione al problema}
Dati tre vettori \(p\), \(q\) e \(r\) rispettivamente di dimensione \(n \times 1\) (colonna), \(1 \times n\) (riga) e
\(n \times 1\) (colonna), si vuole trovare una parentesizzazione del loro prodotto per risolverlo in maniera efficiente.
\[(p \cdot q) \cdot r = (\text{colonna} \cdot \text{riga}) \cdot \text{colonna} = \text{matrice} \cdot \text{colonna} \quad \rightarrow \quad T(n) = \underbrace{\; n^2 \;}_{p_\text{colonna} \; \cdot \; q_\text{riga}} + \underbrace{\; 2n-1 \;}_{\text{matrice} \; \cdot \; r_\text{colonna}}\]
\[p \cdot (q \cdot r) = \text{colonna} \cdot (\text{riga} \cdot \text{colonna}) = \text{colonna} \cdot \text{scalare} \quad \rightarrow \quad T(n) = \underbrace{\; n \;}_{p_\text{colonna} \; \cdot \; \text{scalare}} + \underbrace{\; 2n-1 \;}_{q_\text{riga} \; \cdot \; r_\text{colonna}}\]
Si osserva che parentesizzando in modo diverso, si ottengono complessità decisamente diverse. Elencare tutte le possibili
parentesizzazioni richiederebbe costi esponenziali, per cui è necessario identificare un approccio più efficiente.

\subsubsection*{Formalizzazione del problema}
L'input del problema è una sequenza di \(n+1\) interi \(\langle  p_0, p_1, \dots, p_{n-1}, p_n \rangle\) che rappresentano
le dimensioni delle \(n\) matrici da moltiplicare, secondo la relazione \(A_i \in \mathbb{C}^{p_{i-1} \times p_i}\)
con \(i = 1, \dots, n\).

Inoltre si rappresenta una parentesizzazione di una sequenza di matrici \(\langle  A_i, A_{i+1}, \dots, A_j \rangle\) attraverso
la coppia di indici \(i\) e \(j\) che identificano la prima e l'ultima matrice della sottostringa.

\subsubsection*{Idea risolutiva D\&C e presenza di istanze sovrapposte}
Data una stringa di \(j-i+1\) matrici \(\langle A_i, A_{i+1}, \dots, A_j \rangle\) delimitate dagli indici \(i\) e \(j\),
si identifica la posizione \(k\) ottimale in cui dividere o parentesizzare tale stringa generando due sottostringhe
\(\langle A_i, A_{i+1}, \dots, A_k \rangle\) e \(\langle A_{k+1}, A_{k+2}, \dots, A_j \rangle\) delimitate rispettivamente dalle
coppie di indici \((i, k)\) e \((k+1, j)\). Una volta fatto ciò, si calcola ricorsivamente il costo dei prodotti
necessari a risolvere le due sottostringhe e ottenere le due matrici \(A_{i:k}\) e \(A_{k+1:j}\) a cui poi si aggiungerà
il costo necessario a moltiplicarle tra loro per ottenere la matrice \(A_{i:j}\) corrispondente al prodotto di tutte
le matrici della stringa originale.

Si definisce la funzione \(c[i, j]\) come il costo minimo della moltiplicazione delle matrici della sottostringa
\(\langle A_i, A_{i+1}, \dots, A_j \rangle\) supponendo di aver scelto il valore \(k\) ottimo per minimizzare tale costo.
È utile notare che il costo del prodotto tra le due matrici \(A_{i:k} \in \mathbb{C}^{p_{i-1} \times p_k}\) e
\(B_{k+1:j} \in \mathbb{C}^{p_k \times p_j}\) è dato da \(p_{i-1} \cdot p_k \cdot p_j\).

\[c[i, j] = \begin{cases}
	\; 0 & \text{se } i = j \\
	\; \displaystyle \min_{i \leq k < j} \left\{ \; c[i, k] \; + \; c[k+1, j] \; + \; p_{i-1} \cdot p_k \cdot p_j \; \right\} & \text{se } i < j
\end{cases} \]

\noindent
Si osserva, però, che alcune sottostringhe di matrici possono comparire più volte nell'albero delle chiamate ricorsive,
per cui un approccio puramente ricorsivo con D\&C risulterebbe inefficiente. Si preferisce quindi sfruttare il paradigma
della programmazione dinamica.

\subsubsection*{Soluzione tramite memoizzazione}
Come per il problema LCS, si memorizzano le soluzioni delle sottoistanze identificate dalle coppie di indici \((i, j)\)
all'interno delle celle di una tabella bidimensionale di dimensione \(n \times n\). Si osserva che la tabella ha
le seguenti caratteristiche:
\begin{itemize}
	\item siccome \(i \leq j\), le uniche celle della tabella che vengono effettivamente utilizzate sono quelle sopra
	la diagonale principale, per cui la tabella è triangolare superiore
	\item le celle lungo la diagonale principale \(c[i, j]\) con \(i=j\) contengono il costo \(0\), ovvero il caso base
	\item le celle appena sopra alla diagonale principale \(c[i, j]\) con \(j-i = 1\) riflettono tutte le sottostringhe
	di due matrici, le celle appena sopra a queste riflettono tutte le sottostringhe di tre matrici e così via
	\item la cella \(c[1, n]\) in alto a destra contiene la soluzione finale, ovvero il costo minimo del prodotto di tutte
	le matrici della stringa in input
\end{itemize}

\subsubsection*{Implementazione iterativa}
\vspace{-7pt}
\begin{algo}{IT\_MCM}{$\underline{p}$}{vettore $\underline{p}$ di dimensione $n+1$ con le dimensioni delle matrici}{$m[1, n]$ costo minimo come definito da $c[1,n]$, matrice $\underline{s}$ con i valori di $k$ ottimali}
	\State \(n \gets \text{len}(p) - 1\)
	\For{\(i \gets 1\)}{\(n\)} \(m[i, i] \gets 0\) \Comment{inizializzazione caso base della diagonale principale} \EndFor
	\For{\(l \gets 2\)}{\(n\)} \Comment{per ogni lunghezza \(l\) della sottostringa di matrici}
		\For{\(i \gets 1\)}{\(n-l+1\)} \Comment{discesa lungo la diagonale delle sottostringhe di lunghezza \(l\)}
			\State \(j \gets i + l - 1\)
			\State \(m[i, j] \gets +\infty\) \Comment{inizializzazione con valore sentinella}
			\For{\(k \gets i\)}{\(j-1\)} \Comment{per ogni possibile punto di divisione \(k\)}
				\State \(q \gets m[i, k] + m[k+1, j] + p_{i-1} \cdot p_k \cdot p_j\) \Comment{calcolo del costo per la divisione in \(k\)}
				\If{\(q < m[i, j]\)} \Comment{se il costo calcolato è minore del costo attuale memorizzato}
					\State \(m[i, j] \gets q\) \Comment{aggiornamento del costo minimo per la sottostringa \(i, j\)}
					\State \(s[i, j] \gets k\) \Comment{memorizzazione del punto di divisione ottimo per la sottostringa \(i, j\)}
				\EndIf
			\EndFor
		\EndFor
	\EndFor
	\State \Return \((m[1, n], \underline{s})\)
\end{algo}

\subsubsection*{Alcune osservazioni sull'implementazione iterativa}
\begin{itemize}
	\item il popolamento della tabella avviene per diagonali, partendo da quella principale e procedendo verso l'alto a destra:
	il primo ``for'' varia la lunghezza \(l\) (sceglie la diagonale), il secondo ``for'' varia l'indice \(i\) (sceglie la riga
	della cella da popolare) e l'indice \(j\) viene calcolato di conseguenza in base a \(i\) e \(l\)
	\item in presenza di parentesizzazioni equivalenti (ad esempio per matrici quadrate), l'algoritmo individua quella con
	valore di \(k\) minimo, in quanto utilizzando l'operatore \(<\) e non \(\leq\) nell'if a riga 9, le soluzioni equivalenti
	non vengono considerate migliorative dell'ottimo localmente memorizzato
	\item al posto di avere due tabelle triangolari superiori \(m\) e \(s\), è possibile utilizzare la stessa tabella \(m\)
	popolata superiormente con i costi e inferiormente con i valori di \(k\) ottimali
\end{itemize}

\subsubsection*{Complessità spaziale e temporale}
La complessità temporale dell'algoritmo MCM è \(O(n^3)\) in quanto si visitano circa la metà delle celle della tabella che
asintoticamente sono \(O(n^2)\) e per ciascuna di esse si esegue un ciclo su \(k\) che varia da \(i\) a \(j\) con
\(i, j \in [1, n]\) in cui si svolgono operazioni in tempo costante.

La complessità spaziale dell'algoritmo MCM è \(O(n^2)\) in quanto le tabelle impiegate devono avere dimensione \(n \times n\)
per memorizzare i costi e le posizioni di parentesizzazione ottimali per tutte le possibili coppie di indici \(i\) e \(j\).
A differenza dell'algoritmo di LCS, il valore di una cella \(m[i, j]\) non dipende esclusivamente dalle celle adiacenti, per
cui non si possono fare ottimizzazioni a livello spaziale.

\newpage

\subsection{Longest Palindrome Subsequence - LPS}
\subsubsection*{Definizione di sequenza palindroma}
Una sequenza \(Z = \; \langle z_1, z_2, \dots, z_n \rangle\) si dice palindroma se \(z_{1+h} = z_{n-h}\) per ogni \(h = 0, 1, \dots, n\!-\!1\).

\noindent
Si ricorda che una sottosequenza non necessariamente deve contenere simboli contigui.

\subsubsection*{Definizione della proprietà di sottostruttura}
Per definire la più lunga sottosequenza palindroma di una sequenza \(Z_{i:j}\), si sceglie di confrontare i due simboli
agli estremi \(z_i\) e \(z_j\):
\begin{itemize}
	\item se sono uguali, allora la sottosequenza palindroma conterrà tali due simboli oltre alla più lunga sottosequenza
	palindroma della sottostringa \(Z_{i+1:j-1}\)
	\item se sono diversi, allora la sottosequenza palindroma non conterrà almeno uno dei due simboli, per cui si sceglierà
	la più lunga tra quella della sottostringa \(Z_{i+1:j}\) e quella della sottostringa \(Z_{i:j-1}\)
\end{itemize}
Di seguito si formalizza la proprietà di sottostruttura in termini di equazioni alle ricorrenze.
\[LPS(Z_{i:j}) = \begin{cases}
	\; \left\{ z_i \right\} & \text{se } i = j \\[3pt]
	\; \left\{ z_i, z_j \right\} & \text{se } j - i = 1 \text{ e } z_i = z_j \\[3pt]
	\; \left\{ z_i \right\} \text{ oppure } \left\{ z_j \right\} & \text{se } j - i = 1 \text{ e } z_i \neq z_j \\[3pt]
	\; \left\{ z_i \right\} \;\; \cup \;\; LPS(Z_{i+1:j-1}) \;\; \cup \;\; \left\{ z_j \right\} & \text{se } j-i > 1 \text{ e } z_i = z_j \\[3pt]
	\; \displaystyle \argmax_{\substack{LPS(Z_{i+1:j}) \\ LPS(Z_{i:j-1})}} \left\{ \; \text{len}(LPS(Z_{i+1:j})), \; \text{len}(LPS(Z_{i:j-1})) \; \right\} & \text{se } j-i > 1 \text{ e } z_i \neq z_j
\end{cases}\]
In alternativa, si può definire la proprietà di sottostruttura in termini di lunghezza della sottosequenza palindroma,
nel caso in cui si è interessati solo alla sua lunghezza e non ai simboli che la compongono.
\[p(i, j) = \begin{cases}
	\; 1 & \text{se } i = j \\
	\; 2 & \text{se } j - i = 1 \text{ e } z_i = z_j \\
	\; 1 & \text{se } j - i = 1 \text{ e } z_i \neq z_j \\
	\; p(i+1, j-1) + 2 & \text{se } j-i > 1 \text{ e } z_i = z_j \\
	\; \displaystyle \max \left\{ p(i+1, j), p(i, j-1) \right\} & \text{se } j-i > 1 \text{ e } z_i \neq z_j
\end{cases}\]

\subsubsection*{Note sull'implementazione}
L'implementazione iterativa bottom-up a pagina seguente è molto simile a quella dell'algoritmo per la Longest Common
Subsequence (LCS) tra due stringhe. Come per l'LCS, si utilizza una tabella \(p\) per memorizzare la lunghezza della
sottosequenza palindroma più lunga per ogni sottostringa \(Z_{i:j}\) e una tabella \(b\) per memorizzare a che cella
adiacente è associata alla sottoistanza generata. La principale differenza è che la tabella è riempita a diagonali
come nella Matrix Chain Multiplication (MCM).

Una volta completato l'algoritmo, per ricostruire la sottosequenza palindroma più lunga, si parte dalla cella \(p[1, n]\)
che contiene la lunghezza della sottosequenza palindroma più lunga per l'intera stringa \(Z\) e si procede scendendo
nella tabella \(p\) seguendo le direzioni memorizzate nella tabella \(b\) fino ad arrivare al caso base nelle celle
della diagonale \(p[i, i]\).

\subsubsection*{Complessità computazionale e spaziale}
La complessità computazionale è \(T(n) = \Theta(n^2)\) perché si devono visitare tutte le celle (o al più metà) della
tabella \(p\) di dimensione \(n \times n\) e per ogni cella si effettuano un numero costante di operazioni.

La complessità spaziale è \(S(n) = \Theta(n^2)\) per la tabella \(p\) di dimensione \(n \times n\) e per la tabella \(b\)
di dimensioni analoghe. Come per la LCS, se si è interessati solo alla lunghezza della sottosequenza palindroma più lunga
è possibile ridurre la complessità spaziale a \(S(n) = \Theta(n)\) memorizzando solo le due diagonali della tabella \(p\)
necessarie per soddisfare le dipendenze del calcolo iterativo.

\newpage
\subsubsection*{Implementazione iterativa}
\vspace{-7pt}
\begin{algo}{IT\_LPS}{$\underline{x}$}{sequenza $\underline{x}$}{lunghezza della LPS di $\underline{x}$ e tabella delle transizioni $\underline{b}$ per ricostruire la LPS}
	\State \(n \gets \text{len}(\underline{x})\)
	\For{\(i \gets 1\)}{\(n-1\)} \Comment{inizializzazione dei casi base per le sottosequenze di lunghezza 1 e 2}
		\State \(p[i, i] \gets 1\)
		\If{\(x[i] = x[i+1]\)} \(p[i, i+1] \gets 2\)
		\Else ~ \(p[i, i+1] \gets 1\) \EndIf
	\EndFor
	\State \(p[n,n] \gets 1\)
	\State
	\For{\(l \gets 3\)}{\(n\)} \Comment{calcolo iterativo delle sottosequenze di lunghezza crescente}
		\For{\(i \gets 1\)}{\(n-l+1\)}
			\State \(j \gets i + l - 1\)
			\If{\(x[i] = x[j]\)}
				\State \(p[i, j] \gets p[i+1, j-1] + 2\)
				\State \(b[i, j] \gets \text{``} \swarrow \text{''}\)
			\ElsIf{\(p[i+1, j] \geq p[i, j-1]\)}
				\State \(p[i, j] \gets p[i+1, j]\)
				\State \(b[i, j] \gets \text{``} \downarrow \text{''}\)
			\Else \(\quad (\text{ovvero per } p[i+1, j] < p[i, j-1])\)
				\State \(p[i, j] \gets p[i, j-1]\)
				\State \(b[i, j] \gets \text{``} \leftarrow \text{''}\)
			\EndIf
		\EndFor
	\EndFor
	\State \Return \(p[1, n], \;\; \underline{b}\)
\end{algo}

\begin{algo}{PRINT\_LPS}{$\underline{x}$, $\underline{b}$, $i$, $j$}{sequenza $\underline{x}$, tabella delle transizioni $\underline{b}$, indici $i$ e $j$}{nessun return, stampa la LPS($\underline{x}$) a standard output}
	\If{\(i = j\)} print(\(x[i]\)) \Comment{caso base: sottosequenza di lunghezza 1} \EndIf
	\If{\(j - i = 1\)} \Comment{caso base: sottosequenza di lunghezza 2}
		\If{\(x[i] = x[j]\)} print(\(x[i], x[j]\))
		\Else ~ print(\(x[i]\)) \EndIf
	\EndIf
	\State
	\If{\(b[i, j] = \text{``}\swarrow\text{''}\)} \Comment{simboli agli estremi uguali}
		\State print(\(x[i]\))
		\State \text{PRINT\_LPS}(\(\underline{x}, \underline{b}, i+1, j-1\))
		\State print(\(x[j]\))
	\ElsIf{\(b[i, j] = \text{``}\downarrow\text{''}\)} \Comment{caso 1 - cella sotto}
		\State \text{PRINT\_LPS}(\(\underline{x}, \underline{b}, i+1, j\))
	\Else \(\quad (\text{ovvero per }b[i, j] = \text{``}\leftarrow\text{''}\)) \Comment{caso 2 - cella a sinistra}
		\State \text{PRINT\_LPS}(\(\underline{x}, \underline{b}, i, j-1\))
	\EndIf
\end{algo}
